% -*- Document-Type: LaTeX; Master-File: manual.tex -*- 

% [ck: Aug  9 95]
% \chapter{The standard library}


This chapter presents the main features of the \elan\ standard
library. Remember that the path to this library is specified in the system variable ELANLIB whose default value is set to the subdirectory {\tt elan\_library/elanlib} of the source directory. \elan\ can be used without any reference
to this library, {\it except\/} for what concerns the use of the
built-in objects which are documented in the second part of this chapter.

This   library has  been   designed to  be small  and  as efficient as
possible.   In  particular  {\it  no\/}  AC   operators  is used.  The
resulting  code is more  efficient, at  the  price of sometimes heavier
descriptions.


\section{The \elan\ standard library}

The standard library is fully documented in~\cite{elanlibManual-98} to 
which we refer for details. We give below the list of files defined in
the standard library. Some of them define built-in sorts and operators
and in this case, they also appear in the next section 
(section~\ref{les built-in}).

\subsection{common}

In this category, the most common abstract data types are defined.


\begin{tabular}{l@{~~~~~}l}
        anyIdentifier.eln            &  eq.eln\\
        anyInteger.eln               &  identity.eln\\
        array.eln                    &  int.eln\\
        arrayLIN.eln                 &  integerConstants.eln\\
        arrayLOG.eln                 &  IO.eln\\
        bool.eln                     &  list.eln\\
        builtinArray                 &  occur.eln\\
        builtinHashTree              &  pair.eln\\
        builtinInt.eln               &  prompt.eln\\
        builtinIO.eln                &  replace.eln\\
        builtinSidio.eln             &  stdio.eln\\
        builtinString.eln            &  string.eln\\
        builtinSyntacticMatching.eln &  strlist.eln\\
        cmp.eln                   &  tuple.eln\\
        Query.eln\\
\end{tabular}

\subsection{ref}

All the modules needed to deal with the Reduced \elan\ Format.

\begin{tabular}{l@{~~~~~}l}
        Meta\_Apply.eln         ref\_modules.eln        ref\_terms.eln \\
        REF.eln                ref\_read.eln           ref\_unify.eln \\
        ref\_rules.eln          ref\_unify\_builtin.eln \\
        ref2string.eln         ref\_sorts.eln          ref\_write.eln \\
        ref2term.eln           ref\_strategies.eln     refstring2string.eln \\
        ref\_idents.eln         ref\_tables.eln         string2refterm.eln
\end{tabular}

\subsection{noquote}

These modules describe a possible syntax for the user
specifications. More complicated syntax (e.g. mixfix) can also be
defined.

\begin{tabular}{l@{~~~~~}l}
        applySubstOn.eln   hornClauseSyntax.eln    syntacticUnification.eln \\
        atom.eln           ident.eln               termCommons.eln \\
        atomP.eln          identifier.eln          sigSyntax.eln \\
        eqSystem.eln       substitution.eln
\end{tabular}

\subsection{strategy}

These modules are used for the definition and evaluation of defined
strategies.

\begin{tabular}{l@{~~~~~}l}
        Any.eln          booleg.eln       meta.eln         strsig.eln \\
        any.eln        Meta\_apply.eln   strapply.eln     symbol.eln \\
        Meta\_capply.eln  commit.eln       strat.eln        tcstrat.eln \\
        Meta\_strat.eln   strcapply.eln   loops.eln        strconc.eln
\end{tabular}

\section{The built-ins}\label{les built-in}

%\section{The built-ins}\label{les built-in}
%%% General comment 

There are two types of built-in symbols:
\begin{itemize}
\item 
  {\bf The built-in constructors}, which are non-reducible
  symbols, built into the \elan\ system. Terms built over them are
  either constructed, or interpreted by the \elan\ system, e.g. the
  symbol {\tt Error} defined in the module
  {\ldots/common/builtinIO.eln} as follows:\\
    \begin{verbatim}
    Error(@) : (builtinInt) X   code 122;
    \end{verbatim}
    which is used to report various errors of the I/O sub-system
    written in C++.
\item
    {\bf The built-in functions}, which represent symbols with
    pre-defined semantics, in general, not [easily] specifiable in
    \elan\ itself.  The symbol {\tt open} in  the module
    {\ldots/common/builtinStdio.eln} defined as follows:\\
  \begin{verbatim}
   open_builtinStdio(@,@)     : (builtinString builtinString) builtinInt code 116;
  \end{verbatim}
    can be taken as an example.  In this case, the declaration
    (profile) of the symbol {\tt open} represents an interface between
    the \elan\ specification and its built-in semantics. In this case,
    the code ``116'' identifies its semantics.
    
    If there is an overloaded built-in symbol whose exact semantics
    depends on its profile, a procedure implementing the built-in
    symbol should know, which instance of the built-in symbol is
    executed. The overloaded symbol
    {\tt read\_builtinIO} is an example of this problem:\\
  \begin{verbatim}
    read_builtinIO(@,@)      : (builtinInt X) X   pri 200 code 120;
  \end{verbatim}
    When the built-in symbol {\tt read\_builtinIO} is used for several sorts {\tt
      X}, there exist several built-in symbols {\tt read\_builtinIO} in the
    user's specification.  All of them are linked with the same
    semantic function (with the code 120) parameterized by the
    concrete profile of {\tt read\_builtinIO}. Thus, the sort of a read term is
    known in the semantic procedure due to the profile of actually
    applied symbol {\tt read\_builtinIO}.  The difference between two kinds of
    {\em built-in function} symbols is expressed by the sign of code
    numbers.
\end{itemize}


\subsection{Booleans}

At the beginning there is nothing, so \elan\ provides the true and
false values and introduces the bool module. These two values are
built-in and are deeply connected to the implementation of conditions
in rewrite rules. 

\subsubsection{/elanlib/common/bool.eln}

%% Hubert 

\begin{tabular}{|c|c|c|p{8cm}|}
\hline
{\bf Builtins} & {\bf Code} & {\bf Signature} & {\bf Comments} \\
\hline
\hline
{\tt false}     & 0 & {\tt () bool} & The logical {\tt false}.\\
{\tt true}      & 1 & {\tt () bool} & The logical {\tt true}.\\
{\tt @ and @}   & 21 & {\tt (bool bool) bool} & The logical conjonction.\\
{\tt @ or @}    & 22 & {\tt (bool bool) bool} & The logical disjonction.\\
{\tt not(@)}    & 24 & {\tt (bool) bool} & The logical negation.\\
{\tt @ == @}    & 8 & {\tt (bool bool) bool} & The logical equality test.\\
{\tt @ != @}    & 9 & {\tt (bool bool) bool} & The logical inequality test.\\
{\tt @ $<$ @}   & 10 & {\tt (bool bool) bool} & A logical test assuming that {\tt false<true}.\\
{\tt @ $<$= @}  & 11 & {\tt (bool bool) bool} & A logical test assuming that {\tt false<true}.\\
{\tt @ $>$ @}   & 12 & {\tt (bool bool) bool} & A logical test assuming that {\tt false<true}.\\
{\tt @ $>$= @}  & 13 & {\tt (bool bool) bool} & A logical test assuming that {\tt false<true}.\\
\hline
\end{tabular}


\subsubsection{/elanlib/common/cmp.eln}

%% Hubert

To enrich the booleans, {\it polymorphic\/} equality, disequality and
inequalities are defined and are also built-in:\\

\noindent \begin{tabular}{|c|c|c|p{9cm}|}
\hline
{\bf Builtins} & {\bf Code} & {\bf Signature} & {\bf Comments} \\
\hline
\hline
@ == @          & 18 & {\tt (X X) bool} & Tests the equality between two elements of the same sort {\tt X}.\\
@ != @          & 19 & {\tt (X X) bool} & Tests the inequality between two elements of the same sort {\tt X}.\\
%% @ $<$ @      & 30 & {\tt (X X) bool} & \\
%% @ $<$= @     & 31 & {\tt (X X) bool} & \\
%% @ $>$ @      & 32 & {\tt (X X) bool} & \\
%% @ $>$= @     & 33 & {\tt (X X) bool} & \\
\hline
\end{tabular}



\subsection{Numbers}

Numbers can of course be created ``by hand'', but we choose in \elan\
to provide built-in integers.\\

\subsubsection{/elanlib/common/builtinInt.eln}

%% Hubert

\begin{tabular}{|p{3cm}|c|c|p{4cm}|}
\hline
{\bf Builtins} & {\bf Code} & {\bf Signature} & {\bf Comments} \\
\hline
\hline
{\tt @ + @}                             & 3 & {\tt (builtinInt builtinInt) builtinInt} & The addition on integers.\\
{\tt @ - @}                             & 4 & {\tt (builtinInt builtinInt) builtinInt} & The soustration on integers.\\
{\tt @ * @}                             & 5 & {\tt (builtinInt builtinInt) builtinInt} & The multiplication on integers.\\
{\tt @ / @}                             & 6 & {\tt (builtinInt builtinInt) builtinInt} & The division on integers.\\
{\tt @ \% @}                            & 27 & {\tt (builtinInt builtinInt) builtinInt} & The modulo operation.\\
{\tt @ \& @}                            & 28 & {\tt (builtinInt builtinInt) builtinInt} & {\tt a \& b} makes the logical conjonction between the binary representations of {\tt a} and {\tt b} and returns the corresponding integer.\\
{\tt @ $|$ @}                           & 29 & {\tt (builtinInt builtinInt) builtinInt} & {\tt a $|$ b} makes the logical disjonction between the binary representations of {\tt a} and {\tt b} and returns the corresponding integer.\\
{\tt - @}                               & 20 & {\tt (builtinInt) builtinInt} & The opposite of any integer.\\
{\tt @ == @}                            & 8 & {\tt (builtinInt builtinInt) bool} & The equality test.\\
{\tt @ != @}                            & 9 & {\tt (builtinInt builtinInt) bool} & The inequality test.\\
{\tt @ $<$ @}                           & 10 & {\tt (builtinInt builtinInt) bool} & The strict ordering test between integers.\\
{\tt @ $<$= @}                          & 11 & {\tt (builtinInt builtinInt) bool} & The ordering test between integers.\\
{\tt @ $>$ @}                           & 12 & {\tt (builtinInt builtinInt) bool} & The strict ordering test between integers.\\
{\tt @ $>$= @}                          & 13 & {\tt (builtinInt builtinInt) bool} & The ordering test between integers.\\
{\tt eq\_builtin Int(@,@)}              & 8 & {\tt (builtinInt builtinInt) bool} & The equality test.\\
{\tt neq\_builtin Int(@,@)}             & 9 & {\tt (builtinInt builtinInt) bool} & The inequality test.\\
{\tt less\_builtin Int(@,@)}            & 10 & {\tt (builtinInt builtinInt) bool} & The strict ordering test between integers.\\
{\tt lesseq\_builtin Int(@,@)}          & 11 & {\tt (builtinInt builtinInt) bool} & The ordering test between integers.\\
\hline
\end{tabular}
~\\
\begin{tabular}{|p{3cm}|c|c|p{5,5cm}|}
\hline
{\bf Builtins} & {\bf Code} & {\bf Signature} & {\bf Comments} \\
\hline
\hline
{\tt greater\_builtin Int(@,@)}         & 12 & {\tt (builtinInt builtinInt) bool} & The ordering test between integers.\\
{\tt greatereq\_ builtinInt(@,@)}       & 13 & {\tt (builtinInt builtinInt) bool} & The ordering test between integers.\\
{\tt btoi(@)}                           & 25 & {\tt (bool) builtinInt} & The conversion from a boolean to an integer ($false$ becomes 0 and $true$ becomes 1).\\
{\tt itob(@)}                           & 26 & {\tt (builtinInt) bool} & The conversion from any integer to a boolean (0 becomes $false$ and any other $true$).\\
\hline
\end{tabular}


\subsection{Identifiers}

%Two important built-in modules concern identifiers. First the standard
%ones (i.e. without quotes):\\
%and a similar version but with quotes:\\
%
%In fact quoted identifiers are often used when defining a logic in
%element modules in order to avoid syntactic conflicts at parsing
%time. Unquoted identifiers are mostly used in specifications.\\
%Then we introduce the module {\tt ident} which provides the standard
%operations on identifiers:\\
%and we recommand to use these operators via the module:\\ 
%\insertLib{identifier.eln}\label{identifier module}

\subsubsection{/elanlib/noquote/ident.eln}

%% Hubert

\begin{tabular}{|c|c|c|p{9cm}|}
\hline
{\bf Builtins} & {\bf Code} & {\bf Signature} & {\bf Comments} \\
\hline
\hline
@ == @  & 14 & {\tt (ident ident) bool} & Tests the equality between two identifiers.\\
@ != @  & 15 & {\tt (ident ident) bool} & Tests the inequality between two identifiers.\\
\hline
\end{tabular}



\subsubsection{/elanlib/common/builtinString.eln}

%% Hubert

\begin{tabular}{|c|c|c|p{4cm}|}
\hline
{\bf Builtins} & {\bf Code} & {\bf Signature} & {\bf Comments} \\
\hline
\hline
{\tt strlen(@)}         & 150 & {\tt (string) builtinInt} & Computes the length of a string.\\
{\tt strcat(@,@)}       & 151 & {\tt (string string) string} & Concatenates two strings.\\
{\tt @[@]}              & 152 & {\tt (string builtinInt)}  & {\tt s[i]} selects the {\tt i}th character of the string {\tt s}.\\
{\tt @[@<-@]}           & 153 & {\tt (string builtinInt builtinInt) string} & {\tt s[i$<$-c]} replaces the {\tt i}th character of the string {\tt s} by the character {\tt c}.\\
{\tt substr(@,@,@)}     & 154 & {\tt (string builtinInt builtinInt) string} & {\tt substr(s,i,n)} extracts from the string {\tt s} a substring of length {\tt n} from the {\tt i}th position.\\
{\tt strspn(@,@)}       & 156 & {\tt (string string) builtinInt} & {\tt strspn(s1,s2)} returns the length of the first part of {\tt s1} entirely consituted by characters of {\tt s2}.\\
\hline
\end{tabular}
~\\
\begin{tabular}{|c|c|c|p{5,2cm}|}
\hline
{\bf Builtins} & {\bf Code} & {\bf Signature} & {\bf Comments} \\
\hline
\hline
{\tt strcmp(@,@)}       & 157 & {\tt (string string) builtinInt} & {\tt strcmp(s1,s2)} compares two strings character by character and returns an integer (a negative one if {\tt s1} is less than {\tt s2}, {\tt 0} if they are equal and otherwise a positive integer).\\
{\tt string(@)}         & 158 & {\tt (builtinInt) string} & Conversion from an integer (that is, the associated character) to a string.\\
{\tt ident2string(@)}   & 177 & {\tt (ident) string} & Conversion from an identifier to a string.\\
\hline
\end{tabular}


\subsection{Elementary term computations}

Since it is of primarily use in symbolic computation on terms (and
remember that everything in \elan\ is a term except the built-ins),
the occurrence relation and the replacement operation are provided
as built-ins.

\subsubsection{/elanlib/common/occur.eln}

%% Helene

The module {\tt occur[X,Y]} is parameterized by two sorts {\tt X,Y}
that must be non built-in.\\


\noindent \begin{tabular}{|c|c|c|p{9cm}|}
\hline
{\bf Builtins} & {\bf Code} & {\bf Signature} & {\bf Comments} \\
\hline
\hline
{\tt occurs @ in @}   & 17 & {\tt (X Y) bool} & {\tt occurs s in t} checks if {\tt s} occurs in {\tt t}.\\
\hline
\end{tabular}

\subsubsection{/elanlib/common/replace.eln}

%% Helene

The module {\tt replace[X,Y]} is parameterized by two sorts {\tt X,Y}
that must be non built-in.\\


\noindent \begin{tabular}{|c|c|c|p{8cm}|}
\hline
{\bf Builtins} & {\bf Code} & {\bf Signature} & {\bf Comments} \\
\hline
\hline
{\tt replace @ by @ in @}     & 16 & {\tt (X X Y) Y} & {\tt replace s by u in t} replaces all occurrences of {\tt s} by the {\tt u} in {\tt t}.\\
\hline
\end{tabular}


\subsubsection{/elanlib/common/builtinSyntacticMatching.eln}

%% Helene

Matching and unification operations are provided as built-ins.\\

\noindent \begin{tabular}{|p{4cm}|c|c|p{7,5cm}|}
\hline
{\bf Builtins} & {\bf Code} & {\bf Signature} & {\bf Comments} \\
\hline
\hline
{\tt builtinSyntactic Matching(@,@,@,@)}         & 191 & {\tt (X X Y Y) Y} & {\tt builtinSyntacticMatching (t1,t2,vars,failure)} tries to match terms {\tt t1} to {\tt t2} (i.e. {\tt t1 $<<$ t2}), where {\tt vars:Y} represents list of variables occurred in {\tt t1:X} and {\tt t2:X} and {\tt failure:Y} is a value returned in a non-successful case. Otherwise, a matching substitution is returned in the form of a list of terms associated to the list of variables {\tt vars}. The sort {\tt Y} is supposed to be {\tt list[X]}. It works only for empty theories.\\
\hline
\end{tabular}
~\\
\begin{tabular}{|p{4cm}|c|c|p{7,5cm}|}
\hline
{\bf Builtins} & {\bf Code} & {\bf Signature} & {\bf Comments} \\
\hline
\hline{\tt builtinSyntactic Unification(@,@,@,@)}      & 175 & {\tt (X X Y Y) Y} & the same as the previous built-in function, except that a most general unifier is searched.\\
\hline
\end{tabular}


\subsection{Input/Output}

\subsubsection{/elanlib/common/Query.eln}

%% Helene 

The module {\tt Query[I,O,P]}
defines the parameterized sort Query[I,O,P] and the sort ide 
(embedding the sort ident of identifiers).
All symbols defined in this module are built-in constructors
w.r.t. the classification mentioned before.\\

\noindent \begin{tabular}{|c|c|c|p{6cm}|}
\hline
{\bf Builtins} & {\bf Code} & {\bf Signature} & {\bf Comments} \\
\hline
\hline
{\tt quit}              & 90 & {\tt () Query[I,O,P]} & {\tt quit} just quit the command interpreter\\
{\tt help}              & 112 & {\tt () Query[I,O,P]} & {\tt help} prints the help menu (in the file {\tt elanlib/help.txt}) \\
{\tt load} @            & 91 & {\tt (ide) Query[I,O,P]} & {\tt load f} loads the file {\tt f.eln} and extends the current specification\\
{\tt batch} @           & 107 & {\tt (ide) Query[I,O,P]} & {\tt batch f} runs commands of the batch file {\tt f.eln}\\
{\tt run @}             & 94 & {\tt (I) Query[I,O,P]} & {\tt run q} runs the current specification with the query {\tt q} \\
{\tt sorts @ @ @}       & 92 & {\tt (ide ide ide) [I,O,P]} & {\tt sorts q r s} redefines sorts of query, results and print-outs \\
{\tt startwith (@)@}    & 93 & {\tt (ident O) Query[I,O,P]} & {\tt startwith (s)t} redefines the  starting  term {\tt (s)t} of the {\tt .lgi} file, where {\tt s} is a strategy and {\tt t} is a pattern of the input query\\
{\tt checkwith @}       & 102 & {\tt (bool) Query[I,O,P]} & {\tt checkwith b} redefines the checking term of the {\tt .lgi} file, where {\tt b} is a pattern of a boolean condition tested before any input query\\
{\tt printwith @}       & 98 & {\tt (P) Query[I,O,P]} & {\tt printwith t} redefines the printing term used to pretty-print any result \\
{\tt qs}                & 100 & {\tt Query[I,O,P]} & dump the stack of queries\\
{\tt rs}                & 101 & {\tt Query[I,O,P]} & dump the stack of results\\
{\tt stat}              & 96 & {\tt Query[I,O,P]} & print statistics\\
{\tt dump}              & 95 & {\tt Query[I,O,P]} & dump all rules\\
{\tt dump @ }           & 108  & {\tt (ident) Query[I,O,P]} & {\tt dump(l)} dumps rules or strategies with the label/name {\tt l}\\
{\tt dump @ }           & 109 & {\tt (builtinInt) Query[I,O,P]} & {\tt dump(s)} gives all information about the symbol {\tt s} given by its code\\
{\tt display @}         & 103 & {\tt (builtinInt) Query[I,O,P]} & {\tt display m} changes the printing mode of symbols: {\tt display 1} shows terms in internal form, while {\tt display 0} uses the traditional representation\\
\hline
\end{tabular}
~\\
\noindent \begin{tabular}{|c|c|c|p{6,5cm}|}
\hline
{\bf Builtins} & {\bf Code} & {\bf Signature} & {\bf Comments} \\
\hline
\hline
{\tt trace @}           & 97 & {\tt (builtinInt) Query[I,O,P]} & {\tt trace n} changes the level of debugging to level {\tt n}\\
{\tt break @}           & 104 & {\tt (ident) Query[I,O,P]} & {\tt break f} sets a break-point on rules or strategies  with the label/name {\tt f}\\
{\tt break @}           & 110 & {\tt (builtinInt) Query[I,O,P]} & {\tt break s} sets a break-point on the symbol {\tt s} given by its code\\
{\tt unbreak @}         & 105 & {\tt (ident) Query[I,O,P]} & {\tt unbreak f} deletes the break-point on rules or strategies  with the label/name {\tt f}\\
{\tt unbreak @}         & 111 & {\tt (builtinInt) Query[I,O,P]} & {\tt unbreak s} deletes the break-point on the symbol {\tt s}\\
{\tt breaks}            & 106 & {\tt () Query[I,O,P]} & lists all set break points\\
\hline
\end{tabular}

\subsubsection{/elanlib/common/IO.eln}
%% Huy
The module {\tt IO[X]} defines the Input/Output primitives which can only be used with the compiler.
This module is parameterized by the sort {\tt X} of the terms used in input/output operations.
Currently, only the following sorts are available in the input/output operations: {\tt string}, {\tt int}, {\tt bool}, {\tt term}.
For example, to write out or read in strings, the user only need to import {\tt IO[string]} in his or her program.
%% Hubert

\begin{tabular}{|c|c|c|p{8,5cm}|}
\hline
{\bf Builtins} & {\bf Code} & {\bf Signature} & {\bf Comments} \\
\hline
\hline
{\tt write(@,@)}      & -119 & {\tt (File X) X}          & Writes on output {\tt File} the argument of sort {\tt X}.\\
{\tt writeln(@,@)}      & -119 & {\tt (File X) X}          & Writes on output {\tt File} the argument of sort {\tt X} and goes to the next line.\\
{\tt read(@)}         & -120 & {\tt (File) X}            & Reads from input {\tt File} and return an element of sort {\tt X}.\\
{\tt Error(@)}        & 123  & {\tt (builtinInt) X}     & Reports various errors of the I/O sub-system written in C++.\\
\hline
\end{tabular}

\subsubsection{/elanlib/common/builtinStdio.eln}
%% Huy
The module {\tt builtinStdio.eln} provides several input/output primitives.
But except for {\tt open} and {\tt close}, the others can and should be replaced by the primitives in {\tt IO}.
Notice that the module {\tt builtinStdio.eln} is already imported by {\tt IO} and do not need to be explicitly imported in {\elan} programs.
Three constants of sort {\tt File} {\tt stdin}, {\tt stdout} and {\tt stderr} provide respectively the standard input, the standard output and the standard error output.
%% Hubert

\begin{tabular}{|c|c|c|p{4,5cm}|}
\hline
{\bf Builtins} & {\bf Code} & {\bf Signature} & {\bf Comments} \\
\hline
\hline
{\tt getc(@)}             & 113 & {\tt (builtinInt) builtinInt} & Gets the next character from an input file or pipe.\\
{\tt putc(@)}             & 114 & {\tt (builtinInt builtinInt) builtinInt}  & Puts a character to an output file or pipe.\\
{\tt create(@)}           & 115 & {\tt (builtinString) builtinInt} & Creates a process with a given name. The created process is linked with the parent process via two blocking pipes.\\
{\tt create\_noblock(@)}  & 117 & {\tt (builtinString) builtinInt} & Creates a process with a given name such that two communication pipes are in a non blocking mode.\\
{\tt open(@,@)}           & 116 & {\tt (builtinString builtinString) File} & Opens a file for reading or writing. The first argument is the name of the file, and the second argument is either {\tt "w"} (write or create mode), {\tt "r"} (read mode) or {\tt "a"} (append mode).\\
{\tt close(@)}            & 118 & {\tt (File) builtinInt} & Closes a file or kill a process denoted by the pipe number of sort {\tt File}.\\
\hline
\end{tabular}


\subsection{Strategies}

\subsubsection{/elanlib/strategy/Meta\_strat.eln}

%% Helene

The module {\tt Meta\_strat[X]} parameterized by a sort {\tt X} can
only by used in interpreted mode.  This module exports the sort 
{\tt Strategy[X]}, which specifies an external form of basic strategies
of sort {\tt X}.  The signature of basic strategies is defined by
several built-in constructors used to convert strategy terms of sort
{\tt Strategy[X]} into basic strategies over the sort {\tt X}.  This
conversion is used for the implementation of {\tt meta\_apply} symbols
of the previous module.

~\\
\begin{tabular}{|c|c|p{5cm}|p{6,5cm}|}
\hline
{\bf Builtins} & {\bf Code} & {\bf Signature} & {\bf Comments} \\
\hline
\hline
{\tt @}                 & 131 & {\tt (Strategy[X]) Strategies[X]} & Embeds a basic strategy in a list of basic strategies. \\
{\tt @ , @}             & 132 & {\tt (Strategy[X] Strategies[X]) Strategies[X]} & Concatenates a basic strategy to a list of basic strategies.\\
{\tt @}                 & 133 & {\tt (string) Labels[X]} & Embeds a string in a list of rule labels.\\
{\tt @ , @}             & 134 & {\tt (string Labels[X]) Labels[X]} & Concatenates a string to a list of rule labels.\\
{\tt dc(@)}             & 135 & {\tt (Labels[X]) Strateg[X]} & {\tt dc(l)} is a basic strategy that chooses one label in the list of labels {\tt l} and returns all results of the corresponding rule.\\
{\tt dk(@)}             & 136 & {\tt (Labels[X]) Strateg[X]} & {\tt dk(l)} is a basic strategy that chooses successively all labels in the list of labels {\tt l} and returns all results of the corresponding rules.\\
{\tt dc(@)}             & 137 & {\tt (Strategies[X]) Strateg[X]} & {\tt dc(ls)} is a basic strategy that chooses one basic strategy in the list {\tt ls} and returns all its results.\\
{\tt dk(@)}             & 138 & {\tt (Strategies[X]) Strateg[X]} & {\tt dk(ls)} is a basic strategy that chooses successively all basic strategies in the list {\tt ls}  and returns all their results.\\
{\tt repeat*(@)}        & 139 & {\tt (Strategy[X]) Strateg[X]} & {\tt repeat*(s)} is a basic strategy that repeats the application of the basic strategy {\tt s} until it fails.\\
{\tt iterate*(@)}       & 140 & {\tt (Strategy[X]) Strateg[X]} & {\tt iterate*(s)} is a basic strategy that repeats the application of the basic strategy {\tt s} until it fails and returns intermediate results.\\
{\tt @}                 & 141 & {\tt (Strateg[X]) Strategy[X]} & Embeds the previous constructions into the sort {\tt Strategy[X]}.\\
{\tt @ ; @}             & 142 & {\tt (Strateg[X] Strategy[X])} Strategy[X] & Concatenates any previous construction to a strategy.\\
{\tt id}                & 143 & {\tt () Strateg[X]} & {\tt id} is a basic strategy.\\
{\tt call @ : X}        & 144 & {\tt (string) Strateg[X]} & {\tt call s:X} represents an application of a basic strategy named {\tt s} of sort {\tt X}.
\\ \hline
\end{tabular}


\subsubsection{/elanlib/strategy/Meta\_apply.eln}

%% Helene

The module {\tt Meta\_apply[X]} parameterized by a sort {\tt X} can
only by used in interpreted mode (cf. the module {\tt Meta\_capply[X]}
for the compiled mode).  

The construction {\tt where y:=(META)meta\_apply(s,t)} exported from
this module is semantically equivalent to {\tt where y:=(s)t} where
{\tt s} is a basic strategy, i.e. it gives all results of application
{\tt s} on {\tt t}.  There exists also a sightly modified construction
{\tt where y:=(META)meta\_apply(s,t,n)} giving only the {\tt n}-th
solution of the application of {\tt (s)t}, if it exists.  

The \ELAN\ built-in strategy {\tt META} used here just indicates
that {\tt (META)meta\_apply(s,t)} returns in general several results,
which would not be the case with {\tt ()meta\_apply(s,t)}.  The
built-in strategy {\tt META} is applicable only to terms of the form
{\tt meta\_apply(s,t)} or {\tt meta\_apply(s,t,n)}.

The module {\tt Meta\_apply.eln} exports several variants of the
symbol {\tt set\_of(s,t,m,n)} returning a list of {\tt m} results of
the application {\tt (s)t} from the {\tt n}-th result.  The symbols
{\tt set\_of(s,t)}, {\tt set\_of(s,t,m)} and {\tt set\_of(s,t,m,n)}
are implemented using a result-collecting built-in symbol 
{\tt meta\_apply(s,ts,m,n):(Strategy[X] list[X] builtinInt builtinInt)
  list[X]} defined in this module.  


~\\
\begin{tabular}{|c|c|p{5cm}|p{5cm}|}
\hline
{\bf Builtins} & {\bf Code} & {\bf Signature} & {\bf Comments} \\
\hline
\hline
{\tt meta\_apply(@,@)}          & -129 & {\tt (Strategy[X] X) X} & {\tt meta\_apply(s,t)} returns all results of application of the basic strategy {\tt s} on theterm {\tt t}.\\
{\tt meta\_apply(@,@,@)}        & -130 & {\tt (Strategy[X] X builtinInt) X} & {\tt meta\_apply(s,t,n)} returns the {\tt n}-th result of application of the basic strategy {\tt s} on the term {\tt t}, if it exists.\\
{\tt meta\_apply(@,@,@,@)}      & -128 & {\tt (Strategy[X] list[X] builtinInt builtinInt) list[X]} & The local built-in symbol {\tt meta\_apply(s,t.nil,m,n)} returns {\tt n} results starting from the {\tt m}-th of application of the basic strategy {\tt s} on the term {\tt t}, if they exist.\\
\hline
\end{tabular}


\subsubsection{/elanlib/strategy/Meta\_capply.eln}

%% Helene

The module {\tt Meta\_capply[X]} parameterized by a sort {\tt X} can
only by used in compiled mode.  It represents a simplified form of the
module {\tt Meta\_apply[X]}, where a construction of a strategy of
sort {\tt Strategy[X]} is restricted to a reference of an existing
strategy (i.e. a strategy already occurring in the compiled program).
Thus, in a {\tt meta\_apply} call, only a name (i.e. a string) of a
referenced strategy is used instead of a strategy term of sort 
{\tt   Strategy[X]}.
%%% Its functionality is the same as {\tt Meta\_apply[X]}.

~\\
\begin{tabular}{|c|c|p{4cm}|p{6,5cm}|}
\hline
{\bf Builtins} & {\bf Code} & {\bf Signature} & {\bf Comments} \\
\hline
\hline
{\tt meta\_apply(@,@)}          & -129 & {\tt (string X) X} & {\tt meta\_apply(s,t)} returns all results of application of the strategy named {\tt s} on the term {\tt t}.\\
{\tt meta\_apply(@,@,@)}        & -130 & {\tt (string X builtinInt) X} & {\tt meta\_apply(s,t,n)} returns the {\tt n}-th result of application of the strategy named {\tt s} on the term {\tt t}, if it exists.\\
\hline
\end{tabular}




\subsubsection{/elanlib/strategy/strsig.eln}

%% Helene

The module {\tt strsig[X,Y]} parameterized by the two sorts {\tt X,Y}
defines the signature of the defined strategies.  It introduces several
built-in constructors used to construct strategy terms of the strategy
language.  The module {\tt strsig.eln} contains strategy constructors
used both for type-preserving and type-changing strategies.

~\\
\begin{tabular}{|p{2cm}|c|p{5cm}|p{6cm}|}
\hline
{\bf Builtins} & {\bf Code} & {\bf Signature} & {\bf Comments} \\
\hline
\hline
{\tt [@]@}              & -180 & {\tt ($<$X-$>$Y$>$ X) Y} & {\tt [S]t} applies the strategy {\tt S} to the term {\tt t}.\\
{\tt dk(@)}             & -182 & {\tt (Strlist[X,Y]) $<$X-$>$Y$>$} & {\tt dk(l)} is the defined strategy that chooses successively all strategies in the list {\tt l} and returns all their results.\\
{\tt dc(@)}             & -181 & {\tt (Strlist[X,Y]) $<$X-$>$Y$>$} & {\tt dc(l)} is the defined strategy that chooses one strategy in the list {\tt l} and returns all its results.\\
%one(@)   & -190 & (Strlist[X,Y]) $<$X-$>$Y$>$ & \\
{\tt dc one(@)}         & -190 & {\tt (Strlist[X,Y]) $<$X-$>$Y$>$} & {\tt dc one (l)} is the defined strategy that chooses one strategy in the list {\tt l} and returns one result.\\
{\tt fail}              & -183 & {\tt () $<$X-$>$Y$>$} & The defined strategy that always gives an empty set of results.\\
%%@ @   & -185 & ($<$X-$>$Y$>$ Strlist[X,Y]) Strlist[X,Y] &
%%concatenates a defined strategy to a list of strategies\\
{\tt @,@}               & -185 & {\tt ($<$X-$>$Y$>$ Strlist[X,Y]) Strlist[X,Y]} & Concatenates a defined strategy to a list of strategies.\\
%%@ $\|$ @  & -179 & ($<$X-$>$Y$>$ Strlist[X,Y]) Strlist[X,Y] & \\
%%@ ~,~ @  & -179 & ($<$X-$>$Y$>$ Strlist[X,Y]) Strlist[X,Y] &
%%concatenates a defined strategy to a list of strategies\\
{\tt Epsilon}           & -189 & {\tt () Strlist[X,Y]} & The empty list of strategies.\\
{\tt if @ then @ else @ fi}     & -184 & {\tt (bool $<$X-$>$Y$>$ $<$X-$>$Y$>$) $<$X-$>$Y$>$} & {\tt if b then s1 else s2 fi} is the defined strategy that chooses either {\tt s1} if {\tt b} is true otherwise {\tt s2}.\\
{\tt @}                 & 146 & {\tt (Assignment) AssignmentList} & Embeds an assignment in a list of assignments.\\
{\tt @ , @}             & 147 & {\tt (AssignmentList Assignment) AssignmentList} & Add an assignment to the end of a list of assignments.\\
{\tt if @}              & 145 & {\tt (bool) IfApplication} & A condition if an IfApplication.\\
{\tt @}                 & 146 & {\tt (IfApplication) AssignmentList} & Embeds an IfApplication in a list of assignments.\\
{\tt @ , @}             & 147 & {\tt (AssignmentList IfApplication) AssignmentList} & Add an  IfApplication to the end of a list of assignments.\\
\hline
\end{tabular}


\subsubsection{/elanlib/strategy/strconc.eln}

%% Helene

The module {\tt strconc[X,Y,Z]} parameterized by three sorts
{\tt X,Y,Z} introduces a concatenation symbol $;$ of
two strategies of sorts {\tt <X->Y>} and {\tt <Y->Z>}.

~\\
\begin{tabular}{|c|c|c|p{6cm}|}
\hline
{\bf Builtins} & {\bf Code} & {\bf Signature} & {\bf Comments} \\
\hline
\hline
{\tt @ ; @}     & -188 & {\tt ($<$X-$>$Y$>$ $<$Y-$>$Z$>$) $<$X-$>$Z$>$} & {\tt s1 ; s2} concatenates the strategies {\tt s1:$<$X-$>$Y$>$} and {\tt s2:$<$Y-$>$Z$>$} to get a strategy of sort {\tt $<$X-$>$Z$>$}.\\
\hline
\end{tabular}


\subsubsection{/elanlib/strategy/strat.eln}

%% Helene

The module {\tt strat[X]} parameterized by a sort {\tt X}
has to be imported for working with type-preserving strategies
of sort {\tt <X->X>}.
It defines the {\tt eval} strategy of the interpreter of the language
of defined strategies, and several specific strategy constructors,
which are not present in the case of type-changing strategies.
The rest of strategy constructors is defined in the module {\tt
strategy/strsig.eln}.

~\\
\begin{tabular}{|p{3cm}|c|p{7,5cm}|l|}
\hline
{\bf Builtins} & {\bf Code} & {\bf Signature} & {\bf Comments} \\
\hline
\hline
{\tt id}                        & -186 & {\tt () $<$X-$>$X$>$} & Identity strategy.\\
{\tt if @ then @ orelse @ fi}   & -187 & {\tt ($<$X-$>$X$>$ $<$X-$>$X$>$ $<$X-$>$X$>$) $<$X-$>$X$>$} & Orelse strategy.\\
\hline
\end{tabular}


\subsection{REF Format}

\subsubsection{/elanlib/ref/Meta\_Apply.eln}

%% Helene

The module {\tt ref/Meta\_Apply.eln} provides several variants of the
function {\tt meta\_apply}, which apply a strategy on a term in an
environment defined by a program.  These three components are either
encoded in the REF-format or passed as strings.  Results are produced
equally in the REF-format, or as strings representing terms.

~\\
\begin{tabular}{|c|c|p{5cm}|p{6cm}|}
\hline
{\bf Builtins} & {\bf Code} & {\bf Signature} & {\bf Comments} \\
\hline
\hline
{\tt meta\_apply(@,@,@,@,@)}    & 170 & {\tt (string string string string builtinInt) string} & {\tt meta\_apply(s,t,file,spec,n)} gives the {\tt n}-th result of the application  of the strategy {\tt s} on the term {\tt t} w.r.t. the program stored in files named {\tt file.lgi} and {\tt spec.spc}.\\
{\tt meta\_apply(@,@,@,@,@)}    & 171 & {\tt (string list[string] string string builtinInt) string} & {\tt meta\_apply(s,t.nil,file,spec,n)} gives the {\tt n}-th result of the application of the strategy {\tt s} on the term {\tt t} w.r.t. the program in files named {\tt file.lgi} and {\tt spec.spc}.\\
{\tt meta\_apply(@,@,@,@)}      & 172 & {\tt (string string string builtinInt) string} & {\tt meta\_apply(s,t,file,n)} gives the {\tt n}-th result of the application of the strategy  {\tt s} on the term {\tt t} w.r.t. the program in the REF-format stored in file named {\tt file.ref}.\\
{\tt meta\_apply(@,@,@,@,@)}    & 173 & {\tt (string list[string] string string builtinInt) string} & {\tt meta\_apply(s,t.nil,file,n)} gives the {\tt n}-th result of the application of the strategy {\tt s} on the term {\tt t} w.r.t. the program in the REF-format stored in file named {\tt file.ref}.\\
%meta\_apply(@,@,@,@) & 163 & (string string string builtinInt) string &
%meta\_apply(S,t:sort,P,n-th) gives the n-th result of the application of the
%strategy S on the term t w.r.t. the program P \\
%meta\_apply(@,@,@,@,@) & 164 & (string list[string] string builtinInt builtinInt) list[string] & \\
\hline
\end{tabular}

%\bigskip
%
%\noindent
%Built-in versions of these functions work over strings of REF-terms,
%while the symbols exported from this module work over REF-terms.
%
%~\\
%\begin{tabular}{|c|p{5cm}|p{5cm}|}
%\hline
%{\bf Functions} & {\bf Signature} & {\bf Comments} \\
%\hline
%\hline
%{\tt meta\_apply(@,@,@,@,@)}    & {\tt (Strategy Term Sort Program builtinInt) Term} & {\tt meta\_apply(S,t:sort,P,n)} gives the {\tt n}-th result of a strategy application {\tt S} on the term {\tt t} w.r.t. the program {\tt P}.\\
%{\tt meta\_apply(@,@,@,@)}      & {\tt (string string Program builtinInt) string} & {\tt  meta\_apply(strat,term,P,n)} is like the previous one, while both the strategy {\tt S}, the term {\tt t} and the obtained results are represented in the format REF.\\
%{\tt meta\_apply(@,@,@,@,@,@)}  & {\tt (Strategy Term Sort Program builtinInt builtinInt) list[Term]} & {\tt meta\_apply(S,t:sort,P,n,m)} gives a list of {\tt m} results from the {\tt n}-th one of a strategy application {\tt S} on the term {\tt t} w.r.t. the program {\tt P}. The first result has the index 0.\\
%{\tt meta\_apply(@,@,@,@,@)}    & {\tt (string list[string] Program builtinInt builtinInt) list[string]} & The string version of the previous function.\\
%\hline
%\end{tabular}



%\subsection{/elanlib/ref/divers2string.eln}

%% Helene

%\begin{tabular}{|c|c|c|p{5cm}|}
%\hline
%{\bf Builtins} & {\bf Code} & {\bf Signature} & {\bf Comments} \\
%\hline
%\hline
%type2string(@)  & -176 & (X) string & \\
%ident2string(@) & 177 & (ident) string & \\
%\hline
%\end{tabular}

\subsubsection{/elanlib/ref/ref2string.eln}

%% Helene

The module {\tt ref/ref2string.eln} provides two conversion functions
between strings of REF-terms and terms of the user's defined signature.
It also exports a pretty-printing and a parsing function parameterized
by a signature in the REF-format.

~\\
\begin{tabular}{|l|c|c|p{8cm}|}
\hline
{\bf Builtins} & {\bf Code} & {\bf Signature} & {\bf Comments} \\
\hline
\hline
{\tt refstring2term(@)}         & -168 & {\tt (string) X} & {\tt refstring2term(s)} builds a term of sort {\tt X} from the string {\tt s}.\\
{\tt term2refstring(@)}         & -167 & {\tt (X) string} & {\tt term2refstring(t)} transforms the term {\tt t} of sort {\tt X} into a string.\\
{\tt term2string(@)}            & -169 & {\tt (X) string} & {\tt term2string(t)} pretty-prints the term {\tt t} of sort {\tt X}.\\
\hline
\end{tabular}


\subsubsection{/elanlib/ref/refstring2string.eln}

%% Helene

The module {\tt ref/refstring2string.eln} provides built-in
pretty-printing and parsing functions converting strings of terms into
strings of REF-terms, and vice-versa.
%
This module defines also two symbols {\tt spec2ref} and 
{\tt ThisProgram} constructing a REF-representation of a program.

~\\
\begin{tabular}{|c|c|p{3cm}|p{6,2cm}|}
\hline
{\bf Builtins} & {\bf Code} & {\bf Signature} & {\bf Comments} \\
\hline
\hline
{\tt spec2ref(@,@)}             & 162 & {\tt (string string) string} & {\tt spec2ref(file1,file2)} builds a REF-file from the specification
described by two file {\tt file1.spc} and {\tt file2.lgi}, and it returns its name {\tt file.ref}.\\
{\tt ThisProgram}               & 165 & {\tt () Program} & Produces the REF-format of the running program.\\
{\tt refstring2string(@,@)}     & 160 & {\tt (string string) string} & {\tt refstring2string(s,file.ref)} pretty-prints {\tt s} according to
the grammar defined in the file {\tt file.ref}.\\
{\tt string2refstring(@,@,@)}   & 161 & {\tt (string string string) string} & {\tt string2refstring(st:so,file.ref)} parses the string {\tt st} of a term of sort {\tt so} according to the grammar defined in the file {\tt file.ref}.\\
\hline
\end{tabular}


\section{Calling external programs}

\elan\ allows using external programs like for instance the UNIF
system~\cite{Ringeissen-RTA97}. Let us describe how it works.

The external program is assumed to take as input a term and to return
to \elan\ a set of terms. One could imagine the input term as an input
constraint and the returned terms as the solved forms computed by the
program. The program is called as a sub-process of the \elan\ session
and it is linked to it via two UNIX pipes which are respectively
attached to the standard input and output of the process (i.e. stdin,
stdout). 

Since it is useful to reuse a given process as much as possible,
the description of the called process gives the number of
terms that can be sent to the process before killing it and replacing
it by another one.

Data (i.e. terms) are sent in and out under a textual format. This
format is determined by the signature of the module in which the
process call is performed.

The process call is realized through the use of the strategies
\rw{dccall} and \rw{dkcall} (that stand for {\em dont care
call} and {\em dont know call}), with three arguments:\\
-- the first is the name of the process,\\
-- the second is the number of terms that should be transmitted to the
process before killing it,\\
-- and the third is the (unique) sort of initial and result terms.\\
The syntax is therefore the following:

\grrule{strategy}{
  \rw{dkcall}\sym(\nt{identifier}\sym,\nt{number}\sym,\nt{sort name}\sym)
\gror
  \rw{dccall}\sym(\nt{identifier}\sym,\nt{number}\sym,\nt{sort name}\sym)
}

The difference between the two calls is that the first one returns all
the computed results while the second one returns only the first result,
even if several are computed.

In order to be able to communicate with \elan, the process which is called
should obey to the following syntactic restrictions.\\
The output should have the following syntax:
\grrule{process output}{
  \itone{\sym{\#}} \nt{next solution}
}
\grrule{next solution}{
  \nt{number} \itone{\sym{*}} \nt{term} \itone{\sym{\#}} \nt{next solution}
\gror
  \rw{END} \itone{\sym{\#}}
}
where \nt{term} is one result.\\
In the same way, the process should accept the syntax of the initial
term, described as follows:

\grrule{process input}{
  \nt{term} \nt{next solution demand}
}

\grrule{next solution demand}{
  \sym. \nt{next solution demand}
\gror
  \sym;
}
where \nt{term} is the term passed to the built-in process.

The way the communication is organised between \elan\ and the built-in
process is as follows. \elan\ first sends the initial term to the
process and every time it needs another result, it sends the symbol
``\sym.'' and waits for the answer of the process on its standard
output. The answer should be in the syntax corresponding to the
non-terminal \nt{next solution}. If \elan\ does not need anymore
solution, it sends the symbol ``\sym;'' which makes the process stop,
or it sends an ``end of file'' if the limit of the process use is
reached.

\Attention{Since the synchronization between \elan\ and the called
process in made through pipes, the two process should reinitialise
their buffers after each large output. As a consequence the called
process should call the fflush(stdout) procedure after returning each
solution.} 

\EX

Assume that one would like to use the shell command
``sort'' for sorting a list of ``elements''. The first step is to create a
file containing the shell script calling the sort command and that
follows the syntactic conventions required by \elan\ as just described
above. Let us call this file \texttt{elansort}:
\begin{verbatim}
#! /bin/sh
cat /dev/null > /tmp/inputFile
read r
while (test "$r" != "") do
  echo $r >> /tmp/inputFile 
  read r 
done
echo "##0**"
sort < /tmp/inputFile
echo "##END#"
read r
while (test "$r" != ";") do 
  read r 
done
/bin/rm /tmp/inputFile
\end{verbatim}

Then, we construct the syntax of ``elements'' by using a preprocessor
instruction involving a list (\texttt{Vars}) of identifiers taken from
a specification file:
\begin{verbatim}
import  global  List ;
        local   Vars anyIdentifier list[identifier] ;
end

operators global
    FOR EACH SS:identifier SUCH THAT SS:=(listExtract) elem(Vars) :{
      SS     : elem ;
    }
end
\end{verbatim}

Lists are defined as sequences of elements, without any separator:
\begin{verbatim}
sort list; end

operators global
  @             : ( elem ) list;
  @ @           : ( list list ) list    assocLeft;
end
\end{verbatim}

Since \texttt{elansort} could only be used for sorting elements given
successively on each line of a file, we have to convert lists in the
right syntax before calling the external program:
\insertExamples{listSort.eln}

The main strategy \texttt{extSort} applies first the conversion rule
and then applies the rule calling the external program through the
\texttt{callSort} strategy defined in the following \elan\ module:
\insertExamples{callSort.eln}
This module can be considered as an encapsulation of the
\texttt{elansort} program. According to the alias declarations,
results are written in the initial syntax, as sequences of elements
separated by spaces.
\FEX
